\par\SourceLeaf{525}The rarer this inclination is, the more pleased I am to find you so strongly drawn to the part of mathematics that has always been my favourite study. I sincerely wish you circumstances in which you remain as far as possible master of your time and of the choice of your work. Immediately after the appearance of my Disquisitiones, other kinds of work, and later my circumstances, greatly hindered me from following my inclination as much as I would have wished. Instead of a second part of that work, which I formerly intended, I shall in all probability have to confine myself to presenting a memoir on an individual subject from time to time.
The three memoirs of this kind that have so far appeared in the 16th volume of the local Commentationes and in the first and fourth volumes of the Commentationes recentiores, however, contain (apart from a portion of the second) none of the subjects I already had in mind in 1801 for the continuation, but new ones. My later work of this kind likewise concerns a new subject, namely the theory of biquadratic residues, which I intend to present in about three memoirs; the first of these will shortly be printed for the sixth volume of the Comment. rec., and the main material for the remainder, as well as for the similar theory of cubic residues, may be regarded as essentially settled, although little of it has yet been properly put on paper.
\par Please give my regards to Herr von \textsc{Humboldt}, if he is still in Paris, and excuse my not writing to him directly now by my concern that the letter might not reach him, since, as I hear, he intended to leave Paris.
\par\Needspace{5\baselineskip}With sincere esteem
\par\nobreak\noindent\begin{tabular*}{\linewidth}{@{}l@{\extracolsep{\fill}}r@{}}
&Yours most devotedly\\
Göttingen, 13 September 1826.&C. F. \textsc{Gauss}.\\
\end{tabular*}
\BookRule{.25\textwidth}
\par\endgroup
